% SGA 3, Expose VI_B, Loop 1 English reconstruction.
% Scope: Counterexample 1.2.2 and Proposition 1.3, including
% Lemma 1.3.0 and the complete proof of Proposition 1.3;
% Corollary 1.3.1 is excluded.
% Authority: Polo--Gille Exp6B-13oct24.pdf, component-local pages 3--4
% (combined-reader pages 344--345; original-margin printed page 322,
% ending at the seam before printed page 323).
% OCR and the Reinhold English Markdown are comparison witnesses only.

\medskip
\phantomsection\label{ctr:sga3-VIB-1.2.2}%
\noindent\textbf{\SGARefLink{ctr:sga3-VIB-1.2.2}{反例 1.2.2}.---}
\emph{设 \(k\) 是特征零的域, \(G\) 是常 \(k\)-群 \(\mathbb Z\),
\(H\) 是加法 \(k\)-群 \(\mathbb G_{a,k}\).
设 \(u:G\to H\) 为 \(k\)-群态射.
若 \(u\ne0\), 则 \(u(G)\) 在 \(H\) 中不闭.}

\medskip
\phantomsection\label{prop:sga3-VIB-1.3}%
\noindent\textbf{命题 1.3.---}%
\footnote{\label{note:sga3-VIB-5}编者注: 原文中
\SGARefLink{prop:sga3-VIB-1.3}{1.3},
\SGARefLink{cor:sga3-VIB-1.3.1}{1.3.1} 和
\SGARefLink{lem:sga3-VIB-1.3.1.1}{1.3.1.1} 是在 \(A\) 为域时陈述的.
为求完整, 这里处理 Artin 局部环的情形,
并为此补入了 \hyperref[lem:sga3-VIB-1.3.0]{引理~1.3.0}.}
\emph{设 \(A\) 是剩余域为 \(k\) 的 Artin 局部环,
\(G\) 是局部有限型且平坦的 \(A\)-群,
\(u:G\to H\) 是 \(A\)-群态射. 以下断言等价:
\begin{enumerate}
\item[\textup{(i)}]\phantomsection\label{concl:sga3-VIB-1.3-i}
  \(u\) 在 \(G\) 的某点平坦
  (分别地, 拟有限、非分歧、平滑或平展);%
  \footnote{\label{note:sga3-VIB-6}编者注: 回顾一下,
  概形态射 \(f:X\to Y\) 称为在 \(x\) 处为有限型
  (分别地, 有限呈示或拟有限), 是指存在包含 \(f(x)\) 的仿射开子集
  \(V\) 和包含 \(x\) 的仿射开子集 \(U\), 使 \(f(U)\subset V\),
  且 \(\mathcal O_X(U)\) 是有限型的 \(\mathcal O_Y(V)\)-代数
  (分别地, 为有限呈示, 或还满足对每个 \(x'\in U\),
  \(f^{-1}(f(x'))\) 有限这一附加条件);
  另见第~V~篇的编者注 40.
  若态射在每个点 \(x\) 都有该性质,
  则称其为局部拟有限 (分别地, 局部有限型或局部有限呈示).
  它在 \(x\) 处平滑 (分别地, 非分歧或平展),
  是指存在 \(x\) 的开邻域 \(U\), 使 \(f|_U:U\to Y\)
  平滑 (分别地, 非分歧或平展).
  由这些定义可知, \(f\) 为有限呈示、有限型、拟有限、平滑、非分歧或平展
  的点的集合在 \(X\) 中开. 此外, 由于局部有限呈示态射的平坦点集是开的
  (EGA~IV\(_3\), 11.3.1), \(f\) 同时为有限呈示且平坦的点集也开.
  下文将反复使用这些事实.}
\item[\textup{\SGAInlineRefTarget{prop:sga3-VIB-1.3:item:ii}(ii)}]
  \(u\) 平坦 (分别地, 拟有限、非分歧、平滑或平展).
\end{enumerate}}

\emph{证明.}
只须证明 \textup{\SGARefLink{concl:sga3-VIB-1.3-i}{(i)}} 蕴含
\textup{\SGARefLink{prop:sga3-VIB-1.3:item:ii}{(ii)}}.
先证明下述引理.

\medskip
\phantomsection\label{lem:sga3-VIB-1.3.0}%
\noindent\textbf{引理 1.3.0.---}
\emph{设 \(A\to B\to C\) 为交换环态射,
\(\mathfrak n\) 为 \(A\) 的幂零理想.
假设 \(C/\mathfrak nC\) 是有限型的 \((B/\mathfrak nB)\)-代数.
\begin{enumerate}
\item[\textup{\SGAInlineRefTarget{lem:sga3-VIB-1.3.0:item:i}(i)}]
  则 \(C\) 是有限型的 \(B\)-代数.
\item[\textup{\SGAInlineRefTarget{lem:sga3-VIB-1.3.0:item:ii}(ii)}]
  此外, 若 \(C\) 在 \(A\) 上平坦,
  且 \(C/\mathfrak nC\) 是有限呈示的 \((B/\mathfrak nB)\)-代数,
  则 \(C\) 是有限呈示的 \(B\)-代数.
\end{enumerate}}

事实上, 设 \(x_1,\ldots,x_n\) 是 \(C\) 中的元素,
其像生成作为 \((B/\mathfrak nB)\)-代数的 \(C/\mathfrak nC\).
由幂零形式的 Nakayama 引理, \(x_i\) 生成作为 \(B\)-代数的 \(C\).
这证明了 \SGARefLink{lem:sga3-VIB-1.3.0:item:i}{(i)}.
令
\[
  \varphi:B[X_1,\ldots,X_n]\longrightarrow C
\]
为由此得到的满态射, 并令 \(I=\ker\varphi\).

现假设 \(C\) 在 \(A\) 上平坦, 且 \(C/\mathfrak nC\)
在 \(B/\mathfrak nB\) 上为有限呈示.
一方面, \(I/\mathfrak nI\) 可等同于
\[
  \overline\varphi
    =\varphi\otimes_A(A/\mathfrak n).
\]
的核. 另一方面, 由 EGA~IV\(_1\), 1.4.4,
\(\ker\overline\varphi\) 为有限型理想.
选取 \(P_1,\ldots,P_s\in I\), 使其像生成作为理想的 \(I/\mathfrak nI\).
幂零形式的 Nakayama 引理表明 \(P_i\) 生成 \(I\).
这证明了 \SGARefLink{lem:sga3-VIB-1.3.0:item:ii}{(ii)}.

回到 \hyperref[prop:sga3-VIB-1.3]{1.3} 的证明.
设 \(x\) 为 \(G\) 的任意一点. 由于 \(G\) 在 \(A\) 上平坦,
逐纤维平坦性判别法 EGA~IV\(_3\), 11.3.10.2 表明,
若 \(u\otimes_Ak\) 在 \(x\) 处平坦, 则 \(u\) 也如此.
同样, 前一引理表明, 若 \(u\otimes_Ak\) 在 \(x\) 处为有限型
(分别地, 有限呈示), 则 \(u\) 也如此.
其余性质可在纤维上检验
(非分歧、平滑和平展的情形分别参见 EGA~IV\(_4\),
17.4.1, 17.5.1 和 17.6.1).
因而化归到 \(A=k\) 的情形.

现设 \(x\) 是 \(G\) 中满足 \hyperref[concl:sga3-VIB-1.3-i]{1.3(i)}
的某一条件的点. 所考虑的性质在 \(\mathrm{fpqc}\) 下降下保持
(参见 EGA~IV, 2.5.1, 2.7.1 和 17.7.1).
因此, 将 \(k\) 替换为 \(\kappa(x)\) 的一个代数闭包后,
可化归到 \(k\) 代数闭且 \(x\in G(k)\) 的情形.

由于 \(G\) 是 Jacobson 概形 (参见 EGA~IV\(_3\), 10.4.7),
且 \(G\) 中使 \(u\) 平坦 (分别地, 拟有限、非分歧、平滑或平展)
的点集 \(W\) 在泛化下稳定%
\footnote{\label{note:sga3-VIB-7}编者注: 原文称平坦点集 \(W\)
是开的, 但这一点事先似乎并不显然.}
(分别地, 为开集), 只须证明每个 \(y\in G(k)\) 都属于 \(W\).
对这样的点 \(y\), 平移
\[
  r_y\circ r_x^{-1}
\]
将 \(x\) 映到 \(y\), 因而 \(u\) 在 \(y\) 处有所需性质.
于是 \(y\in W\), 证明完毕.
% Translation statistics: source_lines=132; translated_lines=126; inline_math=108; display_math=3; pending_TODO=0.
% Review: complete component translation; math fragments, reference keys, environments and note counts checked against source; formula prose translated; compilation pending parent integration.
